\chapter{Метастабильность, нуклеация и спинодальный запуск проекторной кристаллизации}

Предыдущая проверка получила точный термический переход
$I_3/3\to\mathbb{RP}^2$, но оставил слово ``охлаждение'' без динамического
содержания. Настоящая проверка сопоставляет найденный функционал с признанной
теорией переходов первого рода и вычисляет всю доступную без новых
коэффициентов структуру метастабильности.

Необходимо различать три задачи:
\begin{enumerate}
 \item почему эффективная обратная температура $\beta$ растёт;
 \item как новая фаза появляется после пересечения равновесного порога;
 \item почему разные области выбирают разные оси и оставляют дефекты.
\end{enumerate}
Литература подробно решает вторую и третью задачи при заданном быстром переключении, но
не заменяет вывода первой задачи из родительской архитектуры.

\section{Точная одномерная фазовая поверхность}

На одноосной семье
\begin{equation}
 R(a)=\operatorname{diag}\left(a,\frac{1-a}{2},\frac{1-a}{2}\right)
\end{equation}
запишем
\begin{equation}
 \mathcal F_\beta(a)=S(a)+\beta E(a),
 \qquad
 S(a)=a\log a+(1-a)\log\frac{1-a}{2},
 \label{eq:v6-projective-kinetic-free-energy}
\end{equation}
где
\begin{align}
 q_2(a)&=\frac{3a^2-2a+1}{2},\\
 q_3(a)&=\frac{3a^3+3a^2-3a+1}{4},\\
 E(a)&=\frac27\left(1-\frac{q_2(a)^2}{q_3(a)}\right)
       +1-q_2(a).
 \label{eq:v6-projective-kinetic-energy}
\end{align}
Это в точности функционал предыдущей проверки после устранения общей
амплитуды моста; новых членов здесь нет.

Ненулевая стационарная ветвь появляется в седло-узловой точке
\begin{equation}
 \boxed{\beta_{\rm ord}=1.3417938971\ldots,\qquad
 a_{\rm ord}=0.7605757614\ldots.}
 \label{eq:v6-ordered-spinodal}
\end{equation}
До неё существует только изотропный минимум. Между
$\beta_{\rm ord}$ и равновесным порогом
\begin{equation}
 \beta_c=1.5426695409\ldots
\end{equation}
одноосная фаза существует, но остаётся метастабильной.

В точке сосуществования между изотропным минимумом и одноосным минимумом
лежит седло
\begin{equation}
 a_\ddagger=0.6178076620\ldots,
 \qquad
 \Delta f_\ddagger=0.0571783574\ldots.
 \label{eq:v6-projective-coexistence-barrier}
\end{equation}
Следовательно, переход первого рода не обязан происходить точно при
$\beta_c$: без флуктуации или зародыша система может продолжить движение
по метастабильной изотропной ветви.

\section{Неожиданно широкое переохлаждение}

Кривизна изотропного состояния вычисляется точно:
\begin{equation}
 \left.\frac{\partial^2\mathcal F_\beta}{\partial a^2}
 \right|_{a=1/3}
 =\frac92-\frac{3\beta}{7}.
 \label{eq:v6-isotropic-curvature-exact}
\end{equation}
Поэтому изотропная фаза теряет локальную устойчивость только при
\begin{equation}
 \boxed{\beta_{\rm sp}=\frac{21}{2}=10.5.}
 \label{eq:v6-isotropic-spinodal-exact}
\end{equation}

Это даёт четыре строгих режима:
\begin{enumerate}
 \item $\beta<\beta_{\rm ord}$: только изотропный минимум;
 \item $\beta_{\rm ord}<\beta<\beta_c$: одноосный минимум метастабилен;
 \item $\beta_c<\beta<\beta_{\rm sp}$: изотропный минимум метастабилен,
 переход требует нуклеации;
 \item $\beta>\beta_{\rm sp}$: изотропный минимум имеет отрицательную
 физическую моду и допускает детерминированный спинодальный распад.
\end{enumerate}

Таким образом, функционал допускает два разных образа рождения мира.
При медленном охлаждении и достаточных флуктуациях возникают отдельные
зародыши. При глубоком переохлаждении гладкая среда долго сохраняется,
после чего теряет устойчивость сразу во всём доступном объёме. Второй
режим ближе к ранней метафоре ``взрыва как свойства самой среды'', но он
пока условен относительно закона $\beta(\tau)$.

\section{Нуклеация: переход через конечный барьер}

Теория Лангера описывает распад метастабильной фазы через критический
зародыш; см. \path{doi:10.1103/PhysRevLett.21.973} и
\path{doi:10.1016/0003-4916(69)90153-5}. Градиентная свободная энергия
неоднородной среды была построена Каном--Хиллиардом; см.
\path{doi:10.1063/1.1744102} и \path{doi:10.1063/1.1730447}.

Если $\sigma$ --- натяжение границы фаз, а $\Delta f$ --- объёмный
выигрыш одноосной фазы, то в трёх измерениях тонкостенная оценка даёт
\begin{align}
 \Delta F(r)&=4\pi\sigma r^2-\frac{4\pi}{3}\Delta f\,r^3,\\
 r_c&=\frac{2\sigma}{\Delta f},\\
 \Delta F_c&=\frac{16\pi\sigma^3}{3\Delta f^2}.
 \label{eq:v6-projective-thin-wall-nucleation}
\end{align}
Проект уже допускает положительную пространственную жёсткость, но не
зафиксировал её абсолютную нормировку. Поэтому форму барьера можно вывести,
а численную скорость нуклеации и радиус пузыря --- ещё нет.

Разные зародыши выбирают независимые элементы
$SO(3)/O(2)=\mathbb{RP}^2$. Их слияние не обязано быть глобально
совместимым; именно на границах доменов возникают проекторные дефекты.
Такой сценарий наблюдается при быстром изотропно-нематическом переходе в
молекулярной динамике; см. \path{doi:10.1103/PhysRevE.65.021705}, а
космологическая аналогия с жидкими кристаллами проверялась в
\path{doi:10.1126/science.251.4999.1336}.

\section{Спинодальный режим: среда сама становится неустойчивой}

Внутри спинодали барьер исчезает: малое возмущение уже уменьшает
свободную энергию. Это стандартное различение нуклеации и спинодального
распада; см. Кан, \path{doi:10.1016/0001-6160(61)90182-1}.

Для ориентационного параметра естественна не сохраняемая релаксационная
динамика типа A,
\begin{equation}
 \partial_\tau R
 =-\Gamma\frac{\delta\mathcal G}{\delta R}+\eta,
 \qquad
 \mathcal G[R]=\int\left(f_\beta(R)+\frac K2|\nabla R|^2\right)d^3x,
 \label{eq:v6-projective-model-a}
\end{equation}
в классификации Хоэнберга--Халперина,
\path{doi:10.1103/RevModPhys.49.435}. При отрицательной локальной кривизне
длинноволновые моды растут, образуя независимо ориентированные области.

Но статический функционал не определяет $\Gamma$, шум $\eta$ и абсолютный
коэффициент $K$. В чистом Model A наиболее быстрая линейная мода находится
при нулевом волновом числе; конечный размер домена появляется из скорости
скорости переключения, причинного горизонта, шума, конечного объёма или дополнительных
сохраняемых полей. Поэтому одна отрицательная кривизна ещё не предсказывает
плотность частиц.

\begin{nogocriterion}[Запрет извлечения плотности дефектов из статики]
Из $\mathcal F_\beta$ и числа $\beta_{\rm sp}$ нельзя вывести размер домена,
скорость роста или число дефектов без кинетического оператора, пространственной
нормировки и закона прохождения через переход.
\end{nogocriterion}

\section{Предел применимости Киббла--Зурека}

Обычный механизм Киббла--Зурека опирается на критическое замедление у
непрерывного перехода; см. \path{doi:10.1088/0305-4470/9/8/029},
\path{doi:10.1038/317505a0} и \path{arXiv:cond-mat/9502119}.
Наш переход строго первого рода: порядок параметра скачет, а при
$\beta_c$ сохраняется конечный барьер \eqref{eq:v6-projective-coexistence-barrier}.
Поэтому стандартную степенную формулу Киббла--Зурека применять напрямую
нельзя.

Современная работа Suzuki--Zurek показывает, что для настраиваемого
слабого перехода первого рода дефектную статистику можно описывать
сочетанием Kibble--Zurek и теории нуклеации; см.
\path{doi:10.1103/PhysRevLett.132.241601}. Для проекта это означает:
топологическая часть Киббла сохраняется, но динамический масштаб должен
приходить из нуклеации или спинодального роста, а не автоматически из
критических показателей непрерывного перехода.

\section{Что литература решила и чего она не решила}

Литературное сопоставление подтверждает следующую цепочку:
\begin{equation}
 \begin{gathered}
  \text{переохлаждение}
  \longrightarrow
  \begin{cases}
   \text{нуклеация пузырей},&\beta_c<\beta<\beta_{\rm sp},\\
   \text{спинодальный распад},&\beta>\beta_{\rm sp},
  \end{cases}\\
  \longrightarrow\mathbb{RP}^2\text{-домены}
  \longrightarrow\text{дефекты}.
 \end{gathered}
 \label{eq:v6-projective-crystallization-chain}
\end{equation}

Однако ни теория Лангера, ни Киббл--Зурек, ни Model A не объясняют, почему
замкнутая безвременная система проекта сама меняет $\beta$. Они принимают
временную эволюцию, тепловой резервуар, расширение или протокол переключения как вход.
Следовательно, механизм фазообразования теперь классифицирован, а механизм
внутреннего охлаждения остаётся отдельным родительским вопросом.

\begin{theorem}[Условная взрывная кристаллизация]
Канонический проекторный функционал обладает широкой метастабильной
изотропной ветвью и точной спинодалью $\beta_{\rm sp}=21/2$. Если внутренняя
динамика доводит систему до этой точки быстрее нуклеации, однородная среда
не может оставаться гладкой и распадается на одноосные $\mathbb{RP}^2$-
домены. Это свойство самой фазовой поверхности, а не внешнего удара.
\end{theorem}

\begin{protocol}
\begin{enumerate}
 \item тип перехода: \textbf{первого рода};
 \item рождение метастабильной одноосной ветви:
 \textbf{$\beta_{\rm ord}=1.3417938971\ldots$};
 \item равновесное сосуществование:
 \textbf{$\beta_c=1.5426695409\ldots$};
 \item барьер в точке сосуществования:
 \textbf{$0.0571783574\ldots$};
 \item точная потеря устойчивости изотропии:
 \textbf{$\beta_{\rm sp}=21/2$};
 \item мелкое переохлаждение: \textbf{нуклеация};
 \item глубокое переохлаждение: \textbf{спинодальный распад};
 \item естественный класс ориентационной динамики:
 \textbf{Model A / Allen--Cahn};
 \item стандартный Kibble--Zurek без модификации:
 \textbf{неприменим};
 \item топологический доменный механизм Киббла:
 \textbf{сохраняется};
 \item скорость, натяжение стенки и плотность дефектов:
 \textbf{не выведены};
 \item закон внутреннего охлаждения $\beta(\tau)$:
 \textbf{не выведен};
 \item следующая проверка:
 \textbf{внутренний перенос энтропии и охлаждение подсистемы};
 \item рождение материи:
 \textbf{кинетически классифицировано, но ещё не замкнуто}.
\end{enumerate}
\end{protocol}

Машинный сертификат сохранён в
\path{s2t_v6_modular_cooling_projective_transition_gate_results.json}.